\documentclass{article}
\usepackage{geometry}
\usepackage{concmath}
\usepackage[OT1]{fontenc}
\usepackage{graphicx}
\usepackage{amsfonts}
\usepackage{amsmath,mathtools}
\usepackage{calrsfs}
\usepackage{color}
\usepackage[dvipsnames]{xcolor}
\usepackage[many]{tcolorbox}
\usepackage{setspace}
\usepackage{tkz-tab}              
\usepackage{multicol}  
\geometry{margin=0.5in}
\definecolor{green}{RGB}{52, 235, 168}
\definecolor{violet}{RGB}{165, 83, 237}
\definecolor{lemon}{RGB}{170, 237, 83}
\definecolor{red}{RGB}{240, 125, 84}
\definecolor{blue}{RGB}{105, 255, 242}

\begin{document}
\noindent \colorbox{green}{220.9}
\begin{equation*}
    A = \begin{pmatrix}
        1&0&1\\-1&2&1\\1&0&1
    \end{pmatrix}
\end{equation*}
a/ \begin{align*}
    &(x,y,z) \in Kerf\\
    &\Leftrightarrow f(x,y,z) = (0,0,0)\\
    &\Leftrightarrow \begin{cases}
        x+y=0\\
        -x+2y+z = 0
    \end{cases}
    \Leftrightarrow \begin{cases}
        z = -x\\
        y = x
    \end{cases}
\end{align*}
Donc, Kerf = Vect(1,1,-1)\\
\begin{align*}
    Imf &= \{(X,Y,Z) | \exists x,y,z, (X,Y,Z) = (x+z,-x+2y+z, x+z)\}\\
    &= Vect(f(e_1), f(e_2), f(e_3))\\
    &= Vect((1,-1,1), (0,2,0))
\end{align*}
* \(Rg(A) = 2\)\\
Car \(C_3 = C_1 + C_2\) et \(C_1\) et \(C_2\) sont libres
* Th du rang : dimKerA = 3-2 = 1\\
* \(C_1 + C_2 - C_3 = 0\)\\
Donc \(\begin{pmatrix}
    1\\1\\-1
\end{pmatrix} \in KerA\)\\
Et donc \(KerA = Vect(\begin{pmatrix}
    1\\1\\-1
\end{pmatrix})\)
* \begin{align*}
    ImA &= Vect(C_1, C_2,C_3)\\
    &= Vect(C_1, C_2)
\end{align*}
\((C_1,C_2,E)\) est une base\\
\begin{equation*}
    \begin{pmatrix}
        1&0&1\\-1&2&1\\1&0&1
    \end{pmatrix}\begin{pmatrix}
        1\\1\\-1
    \end{pmatrix} f:\mathbb{R}^3\to \mathbb{R^3}
\end{equation*}
b/ Soit \(P = Mat(C_1,C_2,E)\)
\begin{equation*}
    det P = -4
\end{equation*}
Donc P est inversible\\
Donc \((C_1,C_2,E)\) est une base de \(M_{3,1}(\mathbb{R})\)\\
Donc \(ImA\oplus KerA = M_{3,1}(\mathbb{R})\)
\begin{align*}
    Mat_{C_1,C_2,E}(f) &= \begin{pmatrix}
        2&0&0\\0&2&0\\0&0&0
    \end{pmatrix}\\
    f(E) &= 0\\
    AC_1 &= \begin{pmatrix}
        1&0&1\\-1&2&1\\1&0&1
    \end{pmatrix}\begin{pmatrix}
        1\\-1\\1
    \end{pmatrix}\\
    &= \begin{pmatrix}
        2\\-2\\2
    \end{pmatrix} = 2C_1 + 0C_2 + 0E\\
    AC_2 &= A\begin{pmatrix}
        0\\2\\0
    \end{pmatrix}\\
    &=\begin{pmatrix}
        0\\4\\0
    \end{pmatrix} = 0C_1 + 2C_2 + 0E
\end{align*}
\begin{equation*}
    f = 2Id \circ p
\end{equation*}

\noindent \colorbox{green}{220.1}\\
\(u(e_1) = e_2\\
u(e_2) = 0\\
u(e_3) = 0\\
e_3 \in Ker u\\
e_2 \in Ker u\)\\
donc \(dim Keru \geq 2\\
e_2 \in Imu\)\\
donc \(Keru \cap Imu \neq \{0\}\)
\(u^2 = 0\) et \(u \neq 0\)\\
donc \(Imu \subset Keru\) et \(Imu \neq \{0\}\)\\
donc \(1 \leq rgu \leq dimKeru\)\\
TDR : rgu + dim Keru = 3\\
donc rg = 1 et dimKeru = 2\\
Soit \(e_1 \notin Ker u\)\\
on pose \(e_2 = u(e_1), e_2 \neq 0\), \(Im u = Vect(e_2) \subset Keru\)\\
on complète \((e_2)\) en \((e_2, e_3)\) base de Keru\\

\noindent \colorbox{green}{220.2}\\
A,B \(\in M_n(\mathbb{R})\)
\(\exists P \in GL_n(\mathbb{R})\) tq \(A = PBP^{-1}\) donc AP = PB\\
où P = Q + iR\\
où \([Q]_{ij} = Re([P]_{ij})\)\\
\([R]_{ij} = Im([P]_{ij})\)\\
On a A(Q+iR) = (Q+iR)B
Donc (AQ-QB0) + i(AR - RB) = 0\\
donc \(\begin{cases}
    AQ-QB = 0\\
    AR-RB = 0
\end{cases}\)\\
Cherchons une CL (réelle) de Q et R inversible en posant 
S(t) = Q+tR\\
On a AS(t) = S(t)B \(\forall t \in R\)\\
J'espère que l'une des S(t) est inversible ! 
\begin{align*}
    f(t) = det(S(t))
    &= \text{fct polynomiale de t}\\
    &= \text{ def du det j'ai la flemme}\\
    \exists ? t_0 tq 
\end{align*}
f(i) = detP \(\neq 0\)\\
donc le polynome f n'est pas nul\\

\noindent \colorbox{green}{220.15}\\
\(C_1 + C_2 - C_3 = 0\)\\
et (\(C_1, C_2\)) libre donc rgA = 2\\
Donc dimKerA = 1\\
Comme \(C_1 + C_2 - C_3 = 0, A\begin{pmatrix}
    1\\1\\-1
\end{pmatrix}=0\)\\
Donc KerA = \(Vect(\begin{pmatrix}
    1\\1\\-1
\end{pmatrix})\)\\
\(A+Id=\begin{pmatrix}
    1&1&1\\1&1&1\\1&1&1 
\end{pmatrix}\)\\
\(C_1 -C_3 = 0\)\\
Donc \(\begin{pmatrix}
    1\\0\\-1
\end{pmatrix} \in Ker(A+Id)\)\\
\(A-Id = \begin{pmatrix}
    -1&1&1\\1&-1&1\\1&1&-1
\end{pmatrix}\)\\
On a \(C_1 + C_2 = 0\)\\
Donc \(\begin{pmatrix}
    1\\1\\0
\end{pmatrix} \in Ker(A-Id)\)\\
Posons \(e_1'\) tq \(Mat_B(e_1')=\begin{pmatrix}
    1\\1\\-1
\end{pmatrix}\)\\
\(e_2'\) ----- \((e_2') = \begin{pmatrix}
    1\\0\\-1
\end{pmatrix}\)\\
\(e_3'\) ----- \((e_3') = \begin{pmatrix}
    1\\1\\0
\end{pmatrix}\)\\
\begin{align*}
    P &= Mat_B(e_1', e_2',e_3') = \begin{pmatrix}
        1&1&1\\1&0&1\\-1&1&0
    \end{pmatrix}\\
    det P &= -1 \neq 0
\end{align*}
Donc \(C = (e_1', e_2',e_3')\) base de E\\
ON a \(Mat_C(f) = \begin{pmatrix}
    0&0&0\\0&-1&0\\0&0&1
\end{pmatrix}\)\\
b/ \begin{equation*}
    A = PDP^{-1}
\end{equation*}\\
c/ \begin{align*}
    X^n &= A^nX_0\\
    &= PD^nP^{-1}\\
    &= \begin{pmatrix}
        1&1&1\\1&0&1\\-1&1&0
    \end{pmatrix}\begin{pmatrix}
        0&0&0\\0&(-1)^n&0\\0&0&1
    \end{pmatrix}\begin{pmatrix}
        -1&-1&-1\\1&-1&0\\1&0&1
    \end{pmatrix}X_0\\
    &= \begin{pmatrix}
        (-1)^n+1\\1\\(-1)^{n+1}
    \end{pmatrix}
\end{align*}

\noindent \colorbox{green}{220.3}\\
\(\forall i \in \{1,\dots,n\}\)\begin{equation*}
    |a_{i,i}|> \sum_{j=1,j\neq i}^{n}|a_{i,j}|
\end{equation*}
Supposons \((C_1,...,C_n)\) liée\\
\(\exists (\lambda_1, ...,\lambda_n) \neq (0,...,0)\)
tq\\
\begin{equation*}
    \sum_{i=1}^{n}\lambda_k C_k = \begin{pmatrix}
        0\\\vdots\\0
    \end{pmatrix}
\end{equation*}
ie \(\begin{cases}
    \sum_{i=1}^{n}\lambda_k a_{1,k} = 0\\
    \vdots\\
    \sum_{i=1}^{n}\lambda_k a_{n,k} = 0
\end{cases}\)\\
Soit \(N \in [1,n]\) tq \(|\lambda_N| \geq |\lambda_i|, \forall i \in [1,n], N\neq i\)\\
On a \begin{align*}
    \sum_{j=1}^{n}\lambda_j a_{N, j}
    \text{ donc } |\lambda_N a_{N,N}| = |\sum_{j=1}^{n} \lambda_j a_{N,j}|\\
     |\lambda_N a_{N,N}| \leq \sum_{j=1}^{n} |\lambda_j a_{N,j}|\\
      |a_{N,N}| < \sum_{j=1}^{n} |\lambda_N a_{N,j}|\\
\end{align*}
Absurde\\
\noindent
\colorbox{green}{220.4}\\
rgA = 1, donc ImA est une droite vectorielle \begin{equation*}
    \exists X \in M_{n,1}(\mathbb{R}), \text{ tq } ImA = Vect(X)
\end{equation*}
* \(C_k \in ImA\) donc \(\exists y_k\) tq  \(C_k = y_k X\)\\
Posons \(Y = \begin{pmatrix}
    y_1\\\vdots\\y_n
\end{pmatrix}\) on a \(A = XY^T\)
b/ \begin{align*}
    A^2 &= XY^TXY^T\\
    &= X(Y^TX)^TY^T\\
    &= \lambda XY^T\\
    &= \lambda A
\end{align*}
Avec \(\lambda = \sum_{k=1}^{n}x_ky_k\)\\

\noindent \colorbox{green}{210.36}\\
a/ \begin{align*}
    f:&E\to E\\
    g:&E \to E\\
    f+g:&E\to E
\end{align*}
Mq Imf+g \(\subset\) Imf + Img\\
Soit \(y \in Imf+g\)\\
Donc \begin{align*}
    \exists x \in E, y &= (f+g)(x)\\
    &= f(x) + g(x)\\
    &\in Imf + Img
\end{align*}
Donc \begin{align*}
    rgf+g &\leq dim(Imf + Img)\\
    &= dimImf + dimImg - dim(Imf\cap Img)\\
    &\leq rgf + rgg
\end{align*}
\begin{align*}
    rgf = rg(fg+(-g))\\
    \leq rg(f+g) + rg(-g)
\end{align*}

\end{document}